% !TEX program = lualatex
\documentclass[11pt]{article}
\usepackage{geometry}
\usepackage{float}
\usepackage{amsmath}
\usepackage{circuitikz}
\usepackage{amssymb}
\usepackage{mathtools}
\usepackage{tikz}
\usepackage{tikz-3dplot}
\usepackage{cite}
\usetikzlibrary{graphs,graphdrawing}  % REQUIRED
\usegdlibrary{layered}
\usetikzlibrary{arrows.meta,graphs,graphdrawing}
\usegdlibrary{layered}
\usepackage{url}
\usetikzlibrary{chains}
\usetikzlibrary{decorations.pathmorphing,calc,positioning,arrows.meta}
\usetikzlibrary{positioning,shapes,arrows.meta}
\usetikzlibrary{circuits.ee.IEC}
\usetikzlibrary{graphs,graphdrawing,quotes}
\usegdlibrary{force}
\usepackage{graphicx}
\usepackage{xcolor}
\usepackage{parskip}
\usepackage{braket}

\title{Obtaining $g$-Factors for Rb-87 and Rb-85 via Optical Pumping}
\author{Spandan Suthar and Antoni Gonzalez}
\date{March 2, 2026}

\geometry{margin=1in}

\begin{document}

\maketitle

\section{Introduction}

% TODO: insert better introduction for optical pumping

Relative to the energy levels described by pure Coulomb-potential
considerations, electrons' energies accrue additional terms in a weak magnetic
field due to several
effects associated with their angular momenta. Three important such
effects are the coupling
between electrons' spins and their own angular momenta,
(spin-orbit, annotated SO), the
coupling between electrons's spins and their nuclei's spins
(hyperfine structure, annotated HS), and the
weak-field (linear) Zeeman effect (annotated Z). Each effect can be treated as a
first-order perturbation to
an electron's original Coulomb Hamiltonian, producing independent
energy contributions:

\begin{equation}
  \braket{\hat H} = \braket{\hat H_\text{C}} + \braket{\hat
  H_\text{SO}} + \braket{\hat H_\text{HS}} +   \braket{\hat H_\text{Z}}
  \label{eq:hamiltonian}
\end{equation}

In a weak magnetic field, electrons' intrinsic spin angular momentum
$\mathbf{S}$ and the electron's
orbital angular momentum $\mathbf{L}$ couple to form a total angular
momentum $\mathbf{J}$,
quantized independently of $\mathbf{S}$ and $\mathbf{L}$. When the
effects of hyperfine structure are considered, $\mathbf{J}$ couples with
the nuclear spin vector $\mathbf{I}$--producing a total atomic angular
momentum $\mathbf{F}$.

When considering spin-orbit
split states subject to hyperfine splitting, electron
states are expressed in a hyperfine-coupled basis, whose members
are eigenstates of the operators $\mathbf{L}^2, \mathbf{S}^2,
\mathbf{J}^2, \mathbf{I}^2, \mathbf{F}^2, F_z$, all simultaneously
diagonalizable. These operators give rise
to the quantum numbers $l, s, j, i, f,$ and $m_f$, respectively, and allow the
states to be labeled as $\ket{n\,\ell\,j\,f\,m_f}$ (omitting $s$ and
$i$, which do not change during transitions).

The perturbative energy shifts added to to the Coulomb Hamiltonian
can each be expressed as a function
of these quantum numbers, and thus of the electrons' states:

\begin{equation}
  \braket{\hat H_\text{SO}}
  = \left\langle
  \frac{\mu_0 e^2}{4\pi m_e^2 r^3}\,
  \hat{\mathbf S}\!\cdot\!\hat{\mathbf L}
  \right\rangle
  = \left\langle
  \frac{\mu_0 e^2}{4\pi m_e^2 r^3}
  \right\rangle
  \frac{\hbar^2}{2}\!\left[
    j(j+1)-\ell(\ell+1)-s(s+1)
  \right]
  \label{eq:spinorbit}
\end{equation}

\begin{equation}
  \braket{\hat H_\text{HS}}
  = \left\langle
  A\,\hat{\mathbf I}\!\cdot\!\hat{\mathbf J}
  \right\rangle
  = A\,\frac{\hbar^2}{2}\!\left[
    f(f+1)-i(i+1)-j(j+1)
  \right]
  \label{eq:hyperfine}
\end{equation}

\begin{equation}
  \braket{\hat H_\text{Z}}
  = \left\langle
  \mu_B\, g_F \,\hat{\mathbf F}\!\cdot\!\mathbf B
  \right\rangle
  = \mu_B\, g_F \, B\,\braket{\hat F_z}
  = \mu_B\, g_F \, B\, \hbar m_f
  \label{eq:zeeman}
\end{equation}

where $A$ is the hyperfine constant, varying from system to system.
Comparing orders of magnitudes, spin-orbit coupling produces the
largest energy shifts---followed by hyperfine structure and the
Zeeman effect.

% Configuration for Rb: [Kr]5s1

This experiment focuses on the $5S_{1/2}$, $5P_{1/2}$, and $5P_{3/2}$ states of
rubidium 85 and rubidium 87 isotopes\footnote{In $nlj$ spectroscopic
  notation, $n$ is the principal quantum number, $l$ is the orbital
  angular momentum, and $j$ is the total electronic angular momentum
  from spin-orbit coupling. $5P_{1/2}$ and $5P_{3/2}$ refer to $n=5$,
  $l=1$ states split into $J=1/2$ and $J=3/2$ fine-structure levels,
respectively.}. Figure \ref{fig:rb87_splitting} depicts the split energy
levels of rubidium 87 due to spin-orbit, hyperfine, and Zeeman
splitting. Transitions between the finest levels can occur for
several reasons, but are always restricted by
selection rules (Table
\ref{tab:e1_selection_rules_corrected}) arising
from the electric-dipole interaction Hamiltonian describing the
interaction of the atom with the electromagnetic field.

\begin{figure}[t]
  \centering
  \resizebox{0.8\textwidth}{!}{\input{rb85_splitting.tikz}}
  \caption{Successive splitting of a representative $^{85}$Rb D1
    excited-state manifold, showing one spin-orbit split ($5P \to
    5P_{1/2},5P_{3/2}$), one hyperfine split ($5P_{1/2}\to F'=2,3$),
    and one Zeeman split ($F'=3\to m_F'$ sublevels).}
  \label{fig:rb85_splitting}
\end{figure}

\begin{figure}[t]
  \centering
  \resizebox{0.8\textwidth}{!}{\begin{tikzpicture}[>=Stealth]

% --- Horizontal positions ---
\def\xA{0}
\def\xAend{2.2}
\def\xB{2.9}
\def\xBend{5.0}
\def\xC{5.7}
\def\xCend{7.8}
\def\xD{8.5}
\def\xDend{11.5}

% --- Vertical positions ---
\def\yP{6}
\def\yS{0}
\def\yPth{7.5}
\def\yPoh{5.2}
\def\ySoh{0}
\def\yFpTwo{5.8}
\def\yFpOne{4.5}
\def\yFgTwo{0.7}
\def\yFgOne{-0.7}
\def\dz{0.22}

% --- Column headers ---
\node[font=\small, anchor=south] at ({(\xA+\xAend)/2}, 8.6) {Coulomb};
\node[font=\small, anchor=south] at ({(\xB+\xBend)/2}, 8.6)
  {$+\hat{H}_{\mathrm{SO}}$};
\node[font=\small, anchor=south] at ({(\xC+\xCend)/2}, 8.6)
  {$+\hat{H}_{\mathrm{HS}}$};
\node[font=\small, anchor=south] at ({(\xD+\xDend)/2}, 8.6)
  {$+\hat{H}_{\mathrm{Z}}$};

% --- Dashed column separators ---
\foreach \xs in {2.55, 5.35, 8.15} {
  \draw[black!25, dashed] (\xs, -1.5) -- (\xs, 8.5);
}

% ===================== UNSPLIT =====================
\draw[very thick] (\xA, \yP) -- (\xAend, \yP);
\node[above, font=\small] at ({(\xA+\xAend)/2}, \yP) {$5P$};

\draw[very thick] (\xA, \yS) -- (\xAend, \yS);
\node[above, font=\small] at ({(\xA+\xAend)/2}, \yS) {$5S$};

% ===================== SO SPLIT =====================
\draw (\xAend, \yP) -- (\xB, \yPth);
\draw (\xAend, \yP) -- (\xB, \yPoh);

\draw[thick] (\xB, \yPth) -- (\xBend, \yPth);
\node[above, font=\small] at ({(\xB+\xBend)/2}, \yPth) {$5P_{3/2}$};

\draw[thick] (\xB, \yPoh) -- (\xBend, \yPoh);
\node[above, font=\small] at ({(\xB+\xBend)/2}, \yPoh) {$5P_{1/2}$};

\draw (\xAend, \yS) -- (\xB, \ySoh);
\draw[thick] (\xB, \ySoh) -- (\xBend, \ySoh);
\node[above, font=\small] at ({(\xB+\xBend)/2}, \ySoh) {$5S_{1/2}$};

% ===================== HS SPLIT =====================
% 5P_{3/2} continues as dashed (D2, not used)
\draw[thick, dashed] (\xBend, \yPth) -- (\xDend, \yPth);

% 5P_{1/2} -> F'=2, F'=1
\draw (\xBend, \yPoh) -- (\xC, \yFpTwo);
\draw (\xBend, \yPoh) -- (\xC, \yFpOne);

\draw[thick] (\xC, \yFpTwo) -- (\xCend, \yFpTwo);
\node[above, font=\small] at ({(\xC+\xCend)/2}, \yFpTwo) {$F'=2$};

\draw[thick] (\xC, \yFpOne) -- (\xCend, \yFpOne);
\node[above, font=\small] at ({(\xC+\xCend)/2}, \yFpOne) {$F'=1$};

% 5S_{1/2} -> F=2, F=1
\draw (\xBend, \ySoh) -- (\xC, \yFgTwo);
\draw (\xBend, \ySoh) -- (\xC, \yFgOne);

\draw[thick] (\xC, \yFgTwo) -- (\xCend, \yFgTwo);
\node[above, font=\small] at ({(\xC+\xCend)/2}, \yFgTwo) {$F=2$};

\draw[thick] (\xC, \yFgOne) -- (\xCend, \yFgOne);
\node[above, font=\small] at ({(\xC+\xCend)/2}, \yFgOne) {$F=1$};

% ===================== ZEEMAN SPLIT =====================
% F'=2: m_F = -2,...,+2
\foreach \mf/\i in {-2/-2, -1/-1, 0/0, 1/1, 2/2} {
  \pgfmathsetmacro\ypos{\yFpTwo + \i*\dz}
  \draw (\xCend, \yFpTwo) -- (\xD, \ypos);
  \draw[thick] (\xD, \ypos) -- (\xDend, \ypos);
  \node[right, font=\scriptsize] at (\xDend, \ypos) {$m_F\!=\!\mf$};
}

% F'=1: m_F = -1,...,+1
\foreach \mf/\i in {-1/-1, 0/0, 1/1} {
  \pgfmathsetmacro\ypos{\yFpOne + \i*\dz}
  \draw (\xCend, \yFpOne) -- (\xD, \ypos);
  \draw[thick] (\xD, \ypos) -- (\xDend, \ypos);
  \node[right, font=\scriptsize] at (\xDend, \ypos) {$m_F\!=\!\mf$};
}

% F=2: m_F = -2,...,+2
\foreach \mf/\i in {-2/-2, -1/-1, 0/0, 1/1, 2/2} {
  \pgfmathsetmacro\ypos{\yFgTwo + \i*\dz}
  \draw (\xCend, \yFgTwo) -- (\xD, \ypos);
  \draw[thick] (\xD, \ypos) -- (\xDend, \ypos);
  \node[right, font=\scriptsize] at (\xDend, \ypos) {$m_F\!=\!\mf$};
}

% F=1: m_F = -1,...,+1
\foreach \mf/\i in {-1/-1, 0/0, 1/1} {
  \pgfmathsetmacro\ypos{\yFgOne + \i*\dz}
  \draw (\xCend, \yFgOne) -- (\xD, \ypos);
  \draw[thick] (\xD, \ypos) -- (\xDend, \ypos);
  \node[right, font=\scriptsize] at (\xDend, \ypos) {$m_F\!=\!\mf$};
}

% --- Energy axis ---
\draw[->, thick] (-0.7, -1.5) -- (-0.7, 8.5);
\node[rotate=90] at (-1.0, 3.5) {Energy};
\node[rotate=90, font=\small\bfseries] at (-1.55, 3.5) {Weak-Field Energy Levels of $^{87}$Rb};

\end{tikzpicture}
}
  \caption{Successive splitting of a representative $^{87}$Rb D1
    excited-state manifold, showing one spin-orbit split ($5P \to
    5P_{1/2},5P_{3/2}$), one hyperfine split ($5P_{1/2}\to F'=1,2$),
    and one Zeeman split ($F'=2\to m_F'$ sublevels). Note that
    splitting due to spin-orbit coupling is much larger than that due
    to hyperfine structure, which is much larger than that due to
  Zeeman splitting.}
  \label{fig:rb87_splitting}
\end{figure}

\begin{table}[H]
  \centering
  \begin{tabular}{cll}
    \hline
    \textbf{} & \textbf{Quantity} & \textbf{Selection Rule (E1)} \\
    \hline
    1 & Spin $s$
    & $\Delta s = 0$ \\

    2 & Orbital $l$
    & $\Delta l = 0,\pm1$ \; ($l=0 \nleftrightarrow l'=0$) \\

    3 & Total electronic $j$
    & $\Delta j = 0,\pm1$ \; ($j=0 \nleftrightarrow j'=0$) \\

    4 & Total atomic $f$
    & $\Delta f = 0,\pm1$ \; ($f=0 \nleftrightarrow f'=0$) \\

    5 & Magnetic sublevel $m_f$
    & $\Delta m_f = 0,\pm1$ \\

    6 & Parity $\pi$
    & $\pi' = -\pi$ \\
    \hline
  \end{tabular}
  \caption{Electric-dipole (E1) selection rules relevant to the $^{87}$Rb D$_1$
    ($5S_{1/2}\rightarrow5P_{1/2}$) optical-pumping transition, expressed in the
  hyperfine-coupled basis.}
  \label{tab:e1_selection_rules_corrected}
\end{table}

% The selection rules and states available produce a nontrivial
% stochastic system; an initial distribution
% of electrons amongst the available states evolves according to
% probabilities of transitions between states, where a transition from state
% $\ket{i}$ to state $\ket{j}$ occurs with probability $P_{ij}$.
%
% \begin{figure}[htbp]
%   \centering
%   \begin{tikzpicture}[
  >=Stealth,
  every node/.style={font=\small, align=center},
  gnode/.style={draw, rounded corners=2pt, minimum width=20mm, minimum height=8mm},
  enode/.style={draw, rounded corners=2pt, minimum width=20mm, minimum height=8mm},
  abs/.style={->, line width=0.7pt},
  em/.style={->, line width=0.6pt},
  rf/.style={->, dashed, line width=0.6pt}
]

\node[anchor=west] at (0,3.2) {\textbf{Edges:} absorption ($\sigma^+$): $\Delta m_F=+1$ \quad emission: $\Delta m_F=0,\pm1$ \quad RF: $\Delta m_F=\pm1$};

\graph [layered layout,
        grow down=16mm,
        level distance=16mm,
        sibling distance=10mm] {

  g_m2[gnode, as={$|g,F=2\rangle$\\$m_F=-2$}];
  g_m1[gnode, as={$|g,F=2\rangle$\\$m_F=-1$}];
  g_0 [gnode, as={$|g,F=2\rangle$\\$m_F=0$}];
  g_p1[gnode, as={$|g,F=2\rangle$\\$m_F=+1$}];
  g_p2[gnode, as={$|g,F=2\rangle$\\$m_F=+2$}];

  e_m1[enode, as={$|e,F'=2\rangle$\\$m_F'=-1$}];
  e_0 [enode, as={$|e,F'=2\rangle$\\$m_F'=0$}];
  e_p1[enode, as={$|e,F'=2\rangle$\\$m_F'=+1$}];
  e_p2[enode, as={$|e,F'=2\rangle$\\$m_F'=+2$}];

  { [same layer] e_m1, e_0, e_p1, e_p2; };
  { [same layer] g_m2, g_m1, g_0, g_p1, g_p2; };

  g_m2 ->[abs] e_m1;
  g_m1 ->[abs] e_0;
  g_0  ->[abs] e_p1;
  g_p1 ->[abs] e_p2;

  e_m1 ->[em] g_m2;
  e_m1 ->[em] g_m1;
  e_m1 ->[em] g_0;

  e_0  ->[em] g_m1;
  e_0  ->[em] g_0;
  e_0  ->[em] g_p1;

  e_p1 ->[em] g_0;
  e_p1 ->[em] g_p1;
  e_p1 ->[em] g_p2;

  e_p2 ->[em] g_p1;
  e_p2 ->[em] g_p2;

  g_m2 ->[rf] g_m1;
  g_m1 ->[rf] g_0;
  g_0  ->[rf] g_p1;
  g_p1 ->[rf] g_p2;

  g_m1 ->[rf] g_m2;
  g_0  ->[rf] g_m1;
  g_p1 ->[rf] g_0;
  g_p2 ->[rf] g_p1;
};

\node[anchor=west] at (10.2,1.2) {\small Excited manifold ($5P_{1/2},F'=2$)};
\node[anchor=west] at (10.2,-1.9) {\small Ground manifold ($5S_{1/2},F=2$)};

\end{tikzpicture}
%   \caption{Five-state transition graph with transition probabilities
%     labeled by $P_{ij}$ for transitions $i\to j$. Nodes represent a subset
%     of $^{87}$Rb D1 Zeeman sublevels in the $f=2$ and $f'=2$ manifolds.
%     Every shown edge is E1-allowed with
%     $\Delta s=0$, $\Delta l=+1$, $\Delta j=0$, and
%   $\Delta m_f\in\{-1,0,+1\}$.}
%   \label{fig:three_state_graph}
% \end{figure}

% In a weak magnetic field, each hyperfine level splits into $m_F$
% sublevels, so the transition energy depends on $\Delta m_F$. The three
% allowed values correspond to the three polarization channels:
% \begin{equation}
%   \Delta m_F = +1 \;(\sigma^+), \qquad
%   \Delta m_F = 0 \;(\pi), \qquad
%   \Delta m_F = -1 \;(\sigma^-).
%   \label{eq:sigma_pi_channels}
% \end{equation}
% In the pumping configuration, circularly polarized light is chosen so
% that repeated absorption favors one sign of $\Delta m_F$, driving
% population toward a stretched state where no further allowed
% $\Delta m_F=+1$ absorption exists.

The rubidium atoms in this experiment are exposed to a magnetic
field $\vec B_0$ along an axis $\vec z$, causing the energy splitting
described in Figure \ref{fig:rb87_splitting} and precession about
$\vec B_0$ with a Larmor frequency $\omega_0$.
The atoms are simultaneously exposed to radiofrequency (RF) pulses
$\vec B_{RF} = B_1 \cos(\omega_1 t) \hat x$ applied transverse to
$\vec B_0$. These RF pulses are indistinguishable from a pair of
counter-rotating magnetic fields of equal strength and frequency
$\omega_{RF}$, circularly polarized in the $x$-$y$ plane:
\begin{align}
  \vec B_{1}(t) &= \frac{B_1}{2}\left[\hat x\cos(\omega_{RF} t)+\hat
  y\sin(\omega_{RF} t)\right] \\
  \vec B_{1}'(t) &= \frac{B_1}{2}\left[\hat x\cos(\omega_{RF} t)-\hat
  y\sin(\omega_{RF} t)\right] \\
  \vec B_{RF} &= \vec B_{1} + \vec B_{1}' = B_1 \cos(\omega_{RF} t) \hat x
\end{align}

Figure \ref{fig:dipole_torque} depicts the geometry of the atom's
total angular momentum $\vec F$ relative to applied fields $\vec B_0$
and $\vec B_1$, including a precession angle $\alpha$ relative to
$\vec B_0$, an azimuthal angular position $\beta$ of $\vec F$, and a
magnetic dipole moment $\vec \mu$ precessing about $\vec F$.

When
$\omega_{RF} = \omega_0 = \dot \beta$,
$\vec B_1$ appears
stationary in the precession frame---while $\vec B_{1}'$ appears to
rotate in the opposite direction.
Both $\vec B_1$ and $\vec B_{1}'$
cause additional torques
on $\vec \mu$. In $\vec \mu$'s frame of reference, $\vec B_1$ appears
stationary while $\vec B_{1}'$ appears to rotate in the opposite
direction with frequency $2\omega_{RF} = 2\omega_0$. Thus, $\vec B_1$ causes
a steady torque $\tau_1$ on $\mu$ while $\vec B_{1}'$ causes a time-dependent
torque $\tau_{1}'$ whose effects net to zero across precession
cycles. The aggregate effect is a steady precession about $\vec B_1$
with a Larmor frequency $\dot \alpha = \omega_1$, causing $\vec F$ to change its
precession angle. Thus, an RF pulse can cause a change in $\vec F_z$,
suggesting that RF pulses can cause changes in the $m_F$ quantum number.

A quantum-mechanical interpretation produces the exact magnitudes of
the RF-driven energy changes. When an RF pulse $\vec B_1(t) = B_1
\cos(\omega_{RF} t) \hat x$ is
applied, the atom's briefly gains a new Hamiltonian contribution due
to the coupling between $\vec \mu$ and $\vec B_1$: $\hat H_{RF} =
-\vec \mu \cdot \left( B_1 \cos(\omega_{RF} t) \hat x \right)$.

As Figure \ref{fig:dipole_torque} illustrates, the atom's magnetic dipole
moment $\vec \mu$ precesses about $\vec F$; in expectation, it is
proportional to $\vec F$: $\braket{\mu} = -\frac{g_F \mu_B}{\hbar}
\vec F$. Thus, one can make the approximation $\vec \mu \approx -\frac{g_F
\mu_B}{\hbar} \vec F$ and $\hat H_{RF}$ becomes:

% H^rf​=−μ⋅Brf​=ℏgF​μB​​F⋅(B1​cos(ωt)x^).

\begin{equation}
  \hat H_{\text{rf}}
  = -\,{\boldsymbol{\mu}}\cdot \mathbf B_{\text{rf}}
  \approx \frac{g_F \mu_B}{\hbar}\mathbf F \cdot \left(B_1 \cos(\omega
  t)\,\hat x\right)
  = \frac{g_F \mu_B B_1}{\hbar}\cos(\omega t)\, \hat F_x
  \label{eq:hamiltonian_rf}
\end{equation}

where $\hat F_x$ is the projection of the total angular momentum
operator $\hat F$ onto the $x$-axis (equivalently, the axis of the RF
pulse). Introduce the operators $\hat F_\pm$ that raise and lower the
$m_F$ quantum number by 1, respectively. These are defined by:
\begin{align}
  \hat F_\pm &= \hat F_x \pm i \hat F_y \quad \left (\text{providing
    } \hat F_x =
  \frac{1}{2} (\hat F_+ + \hat F_-) \right )\\
  \hat F_\pm \ket{F, m_F} &= \hbar \sqrt{F(F+1) - m_F(m_F \pm 1)} \ket{F,
  m_F \pm 1}
  \label{eq:ladder_operators}
\end{align}

Inserting \eqref{eq:ladder_operators} into \eqref{eq:hamiltonian_rf}
yields:

\begin{equation}
  \hat H_{\text{rf}} \ket{F, m_F} \propto (\hat F_+ + \hat F_-) \ket{F,
  m_F} \propto \left( \ket {F, m_F + 1 } + \ket {F,
  m_F - 1 } \right)
\end{equation}

indicating that the RF interaction connects states with relative
$m_F$ values of $\pm 1$, consistent with the classical conclusion
that an RF pulse can change $\vec F_z$ via magnetic torque effects.

\begin{figure}[H]
  \centering
  \begin{tikzpicture}[scale=1.7]
  \definecolor{Fcolor}{HTML}{000000}
  \definecolor{Fproj}{HTML}{505050}
  \definecolor{Mucolor}{HTML}{1F5AA6}
  \definecolor{Bzerocolor}{HTML}{DA4333}
  \definecolor{Bonecolor}{HTML}{FE9920}
  \tdplotsetmaincoords{70}{120}
  \begin{scope}[tdplot_main_coords]
    % Define the total angular-momentum direction (spherical coordinates)
    \def\theta{40}
    \def\phi{150}
    \def\Lradius{2.5}

    % Cartesian coordinates
    \pgfmathsetmacro{\Lx}{\Lradius*sin(\theta)*cos(\phi)}
    \pgfmathsetmacro{\Ly}{\Lradius*sin(\theta)*sin(\phi)}
    \pgfmathsetmacro{\Lz}{\Lradius*cos(\theta)}
    \pgfmathsetmacro{\LxyRadius}{sqrt(\Lx*\Lx + \Ly*\Ly)}

    % Transverse plane
    \filldraw[black!10, opacity=0.3] (-2.5,-2.5,0) -- (2.5,-2.5,0)
    -- (2.5,2.5,0) -- (-2.5,2.5,0) -- cycle;

    % Coordinate axes
    \draw[black,->] (0,0,0) -- (3,0,0) node[anchor=north east]{$x$};
    \draw[black,->] (0,0,0) -- (0,3,0) node[anchor=north west]{$y$};
    \draw[black,->] (0,0,0) -- (0,0,3) node[anchor=south]{$z$};

    % Projection of F onto z
    \draw[thick, dashed, Fproj] (\Lx,\Ly,\Lz) -- (0,0,\Lz);
    \draw[thick, Fproj, ->] (0,0,0) -- (0,0,\Lz);
    \node[Fproj] at (0.5,0,\Lz+0.3) {$\vec{F}_{z}$};

    % Static field B0
    \draw[thick, Bzerocolor, ->] (0,0,0) -- (0,0,1.6);
    \node[Bzerocolor] at (0.5,0,1.0) {$\vec{B}_0$};

    % Precession circle about B0
    \draw[dashed, Bzerocolor] (0,0,\Lz) circle (\LxyRadius);

    % Total angular momentum vector
    \draw[very thick, Fcolor, ->] (0,0,0) -- (\Lx,\Ly,\Lz) node[above right] {$\vec{F}$};

    % Projection of F onto xy
    \draw[thick, dashed, Fproj] (\Lx,\Ly,\Lz) -- (\Lx,\Ly,0);
    \draw[thick, Fproj, ->] (0,0,0) -- (\Lx,\Ly,0);
    \node[Fproj] at (\Lx+0.5,\Ly+0.6,0) {$\vec{F}_{xy}$};

    % Azimuthal angle beta
    \pgfmathsetmacro{\betaAngle}{atan2(\Ly,\Lx)}
    \draw[thick, Bzerocolor] (0,0,0) + (0:0.7) arc (0:{\betaAngle}:0.7);
    \pgfmathsetmacro{\betaLabelAngle}{\betaAngle/2}
    \pgfmathsetmacro{\betaLabelX}{1.1*cos(\betaLabelAngle) - 0.75}
    \pgfmathsetmacro{\betaLabelY}{1.1*sin(\betaLabelAngle) - 0.3}
    \node[Bzerocolor] at (\betaLabelX,\betaLabelY,0) {$\beta$};

    % Rotating field B1 (orthogonal to F_xy in xy-plane)
    \pgfmathsetmacro{\Bnorm}{sqrt(\Lx*\Lx + \Ly*\Ly)}
    \pgfmathsetmacro{\BoneScale}{1.0}
    \pgfmathsetmacro{\BoneX}{\BoneScale*\Ly/\Bnorm}
    \pgfmathsetmacro{\BoneY}{-\BoneScale*\Lx/\Bnorm}
    \draw[thick, Bonecolor, ->] (0,0,0) -- (\BoneX,\BoneY,0);
    \node[Bonecolor] at (\BoneX+0.3,\BoneY+0.3,0) {$\vec{B}_1$};

    % Circle in plane perpendicular to B1 passing through F
    \pgfmathsetmacro{\BoneAngle}{atan2(\BoneY,\BoneX)}
    \pgfmathsetmacro{\BoneNorm}{sqrt(\BoneX*\BoneX + \BoneY*\BoneY)}
    \pgfmathsetmacro{\MdotBone}{(\Lx*\BoneX + \Ly*\BoneY)/\BoneNorm}
    \pgfmathsetmacro{\MparallelX}{\MdotBone*\BoneX/\BoneNorm}
    \pgfmathsetmacro{\MparallelY}{\MdotBone*\BoneY/\BoneNorm}
    \pgfmathsetmacro{\MperpX}{\Lx - \MparallelX}
    \pgfmathsetmacro{\MperpY}{\Ly - \MparallelY}
    \pgfmathsetmacro{\MperpZ}{\Lz}
    \pgfmathsetmacro{\radiusPerpToBone}{sqrt(\MperpX*\MperpX + \MperpY*\MperpY + \MperpZ*\MperpZ)}
    \tdplotsetrotatedcoords{\BoneAngle}{90}{0}
    \tdplotdrawarc[tdplot_rotated_coords,dashed,Bonecolor]{(0,0,0)}{\radiusPerpToBone}{0}{360}{}{}

    % Spin-orbit picture: mu precessing around F
    \pgfmathsetmacro{\ux}{\Lx/\Lradius}
    \pgfmathsetmacro{\uy}{\Ly/\Lradius}
    \pgfmathsetmacro{\uz}{\Lz/\Lradius}
    \pgfmathsetmacro{\rho}{sqrt(\ux*\ux + \uy*\uy)}
    \pgfmathsetmacro{\px}{-\uy/\rho}
    \pgfmathsetmacro{\py}{\ux/\rho}
    \pgfmathsetmacro{\pz}{0}
    \pgfmathsetmacro{\qx}{-\uz*\py}
    \pgfmathsetmacro{\qy}{\uz*\px}
    \pgfmathsetmacro{\qz}{\ux*\py - \uy*\px}

    \pgfmathsetmacro{\muAlong}{1.45}
    \pgfmathsetmacro{\muConeR}{0.55}
    \pgfmathsetmacro{\cx}{\muAlong*\ux}
    \pgfmathsetmacro{\cy}{\muAlong*\uy}
    \pgfmathsetmacro{\cz}{\muAlong*\uz}
    \def\muPhase{270}
    \def\muTilt{14}

    % Slight visual tilt of the mu-precession ring toward the viewer.
    \pgfmathsetmacro{\qtx}{\qx*cos(\muTilt) + \ux*sin(\muTilt)}
    \pgfmathsetmacro{\qty}{\qy*cos(\muTilt) + \uy*sin(\muTilt)}
    \pgfmathsetmacro{\qtz}{\qz*cos(\muTilt) + \uz*sin(\muTilt)}

    \pgfmathsetmacro{\muX}{\cx + \muConeR*(\px*cos(\muPhase) + \qx*sin(\muPhase))}
    \pgfmathsetmacro{\muY}{\cy + \muConeR*(\py*cos(\muPhase) + \qy*sin(\muPhase))}
    \pgfmathsetmacro{\muZ}{\cz + \muConeR*(\pz*cos(\muPhase) + \qz*sin(\muPhase))}

    \draw[dashed, Mucolor]
      plot[samples=80, domain=0:360, smooth, variable=\t]
      ({\cx + \muConeR*(\px*cos(\t) + \qtx*sin(\t))},
       {\cy + \muConeR*(\py*cos(\t) + \qty*sin(\t))},
       {\cz + \muConeR*(\pz*cos(\t) + \qtz*sin(\t))});
    \draw[thick, Mucolor, ->] (0,0,0) -- (\muX,\muY,\muZ)
      node[pos=0.5, below=2pt] {$\vec{\mu}$};

    % Polar angle alpha relative to B0
    \tdplotsetthetaplanecoords{\phi}
    \tdplotdrawarc[tdplot_rotated_coords,thick,Bonecolor]{(0,0,0)}{0.5}{0}{\theta}{}{}
    \pgfmathsetmacro{\labelR}{0.8}
    \pgfmathsetmacro{\labelTheta}{\theta/2}
    \pgfmathsetmacro{\labelX}{\labelR*sin(\labelTheta)*cos(\phi)}
    \pgfmathsetmacro{\labelY}{\labelR*sin(\labelTheta)*sin(\phi)}
    \pgfmathsetmacro{\labelZ}{\labelR*cos(\labelTheta)}
    \node[Bonecolor] at (\labelX,\labelY,\labelZ) {$\alpha$};
  \end{scope}
\end{tikzpicture}

  \caption{Precession geometry showing $\vec{F}$ precessing about a
    magnetic field $\vec{B}_0$ of constant direction. In the
    spin-orbit-coupled picture, the magnetic dipole moment of the
    atom $\vec{\mu}$
    precesses about $\vec{F}$. Note the azimuthal and polar positions
    $\beta$ and $\alpha$ of
    $\vec{F}$, associated with the Larmor frequencies $\omega_0$ and $\omega_1$
    caused by $\vec{B}_0$ and $\vec{B}_1$, respectively.
  }
  \label{fig:dipole_torque}
\end{figure}

Another avenue providing control over the rubidium atoms' energy levels is
interaction with circularly polarized photons in an E1
electric-dipole interaction. During these interactions, E1 selection
rules apply: these include $\pi' = -\pi$, $\Delta \ell = \pm 1$,
$\Delta j = 0, \pm 1$, and $\Delta m_f = q$ determined by the
light's polarization component $q$. Note that $q$ is $+1$ for
right-circularly polarized light and $-1$ for
left-circularly polarized light.

These photons can cause transitions
consistent with the selection rules expressed in Table
\ref{tab:e1_selection_rules_corrected}---namely,
transitions with $\Delta l = \pm 1$, $\Delta j = 0, \pm 1$, and
$\Delta m_f = 0, \pm 1$, allowing transitions between the
ground $5S_{1/2}$ manifold and the excited $5P_{1/2}$ and $5P_{3/2}$
manifolds. Upon spontaneous emission, the atoms return to the ground
$5S_{1/2}$ manifold, resulting in unchanged $l$ and $j$ quantum
numbers. Abstractly, for an atom starting in the ground manifold,
a single absorption-emission cycle can be viewed as a process that changes
$f$, $m_f$, or both.

\[
  \ket{n, l^{(0)}, j^{(0)}, f^{(0)}, m_f^{(0)}}
  \xrightarrow{\text{absorption}}
  \ket{n, l^{(1)}, j^{(1)}, f^{(1)}, m_f^{(1)}}
  \xrightarrow{\text{emission}}
  \ket{n, l^{(0)}, j^{(0)}, f^{(2)}, m_f^{(2)}}
\]
\[
  \Downarrow
\]
\[
  \ket{f, m_f} \xrightarrow{\text{absorption-emission}} \ket{f', m_f'}
\]

Note that other short-lived intermediate states in the excited
manifold can be attained between absorption and emission, not
depicted above for clarity. Also of note is that atoms' stints in the
excited manifold are orders of magnitude more short-lived than atoms
in the ground manifold, suggesting that these intermediate states are
negligible.

Spontaneous emissions have no bias toward a particular $\Delta
m_f$ value. Thus, if the light triggering an absorption-emission cycle is
right-circularly polarized, possible values of $\Delta m_f$ across an
complete absorption-emission cycle are $m_f = \Delta m_f^\text{absorption} +
\Delta m_f^\text{emission} \in \set{-1 + 1, 0 + 1, 1 + 1} = \set{0,
+1, +2}$, creating a bias toward higher $m_f$ values. Consider states
with $F = 2$ and $m_F = 2$. Figure \ref{fig:rb87_splitting} suggests that
a transition with $\Delta m_f = +1$ is forbidden; no states in the
excited manifold have $m_F = 3$. Thus, the stochastic process governing
the redistribution of atomic states through the ground manifold during
repeated absorption-emission cycles ($\mathcal P$) provides no
outward flow from the
state $\ket{\eta} \equiv \ket{n = 5, l = 0, j = 1/2, f = 2, m_f =
2}$\footnote{Allow states to be uniquely indexed as $\ket{i}$ for
  brevity, with $i = \eta$ identifying the unique zero-outflow
state.}---while providing
healthy inward flow. No other states have this property, and $\ket{\eta}$
is reachable from any other state in the ground manifold. Thus,
$\ket{\eta}$ becomes an attractor for the system, with an
attraction basin that is the entire $5S_{1/2}$ ground manifold.

Thus, from a uniform distribution $\pi_0$ of states in the ground manifold,
the distribution $\pi_n$ of states after $n$ absorption-emission
cycles converges to a delta distribution centered on $\ket{\eta}$:

$$
\lim_{n \to \infty} \pi_n = \lim_{n \to \infty} (\mathcal P^n[\pi_0]) =
\delta_{\eta}
$$

The distribution $\delta_\eta$ is characterized by zero photon
absorption probability, causing the rubidium isotope to transmit
light of the wavelength of the $5P_{1/2} \rightarrow 5S_{1/2}$
transition.

$\mathcal P$ can be changed by switching on an RF field $\vec B_1$.
Established above, the RF field produces spontaneous $\Delta m_f = \pm
1$ transitions through coupling with $\vec \mu$, removing the bias toward higher
$m_f$ values and providing outward flow from $\ket{\eta}$. This
produces a new steady-state distribution $\pi_\text{depumped}$,
characterized by a nonzero absorption probability for the $5P_{1/2} \rightarrow
5S_{1/2}$ transition photons and a smaller transmission intensity.

\section{Procedure}

The experimental setup for optical pumping is shown in Figure
\ref{fig:optical_pumping_setup}. Light from an RF discharge lamp
passes through an interference filter to select the appropriate
wavelength for the rubidium D$_1$ transition ($5S_{1/2} \rightarrow
5P_{1/2}$, $\lambda = 794.8$ nm). The light then passes through a
linear polarizer and a quarter-wave plate to produce circularly
polarized light, which is directed through the rubidium absorption
cell. The rubidium cell is placed in a static magnetic field oriented
horizontally and an RF magnetic field oriented vertically. The
transmitted light intensity is measured by an optical detector.

\begin{figure}[H]
  \centering
  \resizebox{\textwidth}{!}{\begin{tikzpicture}[scale=1.2]

  % Centerline arrow
  \draw[-{Stealth[length=2mm,width=1.5mm]},gray!50] (-1.2,0) -- (12.5,0);

  % RF Discharge Lamp
  \begin{scope}[shift={(0,0)}]
    \draw (0,0) circle (0.6);
    \node[align=center,below] at (0,-1.0) {\small RF\\[-0.1cm]\small discharge\\[-0.1cm]\small lamp};
  \end{scope}

  % First Convergent Lens
  \begin{scope}[shift={(1.5,0)}]
    \draw (0,-1) to[bend left=15] (0,1);
    \draw (0,-1) to[bend right=15] (0,1);
    \node[align=center,below] at (0,-1.3) {\small Lens};
  \end{scope}

  % Interference Filter
  \begin{scope}[shift={(3.0,0)}]
    \draw (-0.08,-1) rectangle (0.08,1);
    \node[align=center,below] at (0,-1.3) {\small Interference\\[-0.1cm]\small filter};
  \end{scope}

  % Linear Polarizer
  \begin{scope}[shift={(4.5,0)}]
    \draw (-0.08,-1) rectangle (0.08,1);
    \node[align=center,below] at (0,-1.3) {\small Linear\\[-0.1cm]\small polarizer};
  \end{scope}

  % Quarter Wave Plate
  \begin{scope}[shift={(6.0,0)}]
    \draw (-0.08,-1) rectangle (0.08,1);
    \node[align=center,below] at (0,-1.3) {\small Quarter\\[-0.1cm]\small wave plate};
  \end{scope}

  % Rubidium Absorption Cell (center)
  \begin{scope}[shift={(8.2,0)}]
    \draw (-0.8,-0.5) rectangle (0.8,0.5);
    \node[align=center,below] at (0,-0.85) {\small Rubidium\\[-0.1cm]\small absorption cell};
  \end{scope}

  % Second Convergent Lens
  \begin{scope}[shift={(10.0,0)}]
    \draw (0,-1) to[bend left=15] (0,1);
    \draw (0,-1) to[bend right=15] (0,1);
    \node[align=center,below] at (0,-1.3) {\small Lens};
  \end{scope}

  % Optical Detector
  \begin{scope}[shift={(11.5,0)}]
    \draw (-0.15,-0.4) rectangle (0.15,0.4);
    \node[align=center,below] at (0,-0.8) {\small Optical\\[-0.1cm]\small detector};
  \end{scope}

  % Light rays
  % Upper ray: lamp top → focal point in cell
  \draw (0,0.6) -- (8.2,0);
  % Lower ray: lamp bottom → focal point in cell
  \draw (0,-0.6) -- (8.2,0);

  % Upper ray: focal point → second lens → detector
  % Angle matches: slope_in = -0.6/8.2 ≈ -0.0732, slope_out = +0.0732
  % At x=10.0: y = 0.0732 * 1.8 ≈ 0.132
  \draw (8.2,0) -- (10.0,0.132);
  \draw (10.0,0.132) -- (11.35,0);

  % Lower ray: focal point → second lens → detector
  \draw (8.2,0) -- (10.0,-0.132);
  \draw (10.0,-0.132) -- (11.35,0);

  % Coordinate system indicator (above cell)
  \begin{scope}[shift={(8.2,1.5)}]
    % z-axis (left to right, along optical axis)
    \draw[-{Stealth[length=2mm,width=1.5mm]}] (0,0) -- (0.67,0);
    \node[right] at (0.67,0) {$z$};
    % y-axis (down to up)
    \draw[-{Stealth[length=2mm,width=1.5mm]}] (0,0) -- (0,0.67);
    \node[above] at (0,0.67) {$y$};
    % x-axis (into page)
    \fill[white] (0,0) circle (0.125);
    \node at (0,0) {$\otimes$};
    \node[below left] at (0,0) {$x$};
  \end{scope}

\end{tikzpicture}
}
  \caption{Schematic diagram of the optical pumping apparatus. Light
    from the RF discharge lamp is filtered, polarized, and directed
    through the rubidium absorption cell, which is subjected to both
    static and RF magnetic fields. The transmitted light intensity is
  monitored by the optical detector.}
  \label{fig:optical_pumping_setup}
\end{figure}

\begin{figure}[H]
  \centering
  \resizebox{0.5\textwidth}{!}{\begin{tikzpicture}[scale=1.0]

  % Large Helmholtz coils (annulus)
  \draw (0,0) circle (3.0);
  \draw (0,0) circle (2.8);

  % Top and bottom rectangular coils
  \draw (-1.5,1.2) rectangle (1.5,1.4);
  \draw (-1.5,-1.2) rectangle (1.5,-1.4);

  % Left and right rectangular coils
  \draw (-1.4,-1.0) rectangle (-1.2,1.0);
  \draw (1.2,-1.0) rectangle (1.4,1.0);

  % Rubidium absorption cell at center
  \draw (-0.6,-0.4) rectangle (0.6,0.4);

  % Labels with arrows
  % Outer coil label
  \draw[-{Stealth[length=2mm,width=1.5mm]}] (3.0,3.0) -- (2.1,2.1);
  \node[above right] at (3.0,3.0) {$\mathbf{B}_0$ coils};

  % Vertical coils label
  \draw[-{Stealth[length=2mm,width=1.5mm]}] (2.8,1.3) -- (1.5,1.3);
  \node[right] at (2.8,1.3) {Vertical field control coils};

  % RF coils label
  \draw[-{Stealth[length=2mm,width=1.5mm]}] (-3.2,0) -- (-1.4,0);
  \node[left] at (-3.2,0) {$\mathbf{B}_\text{rf}$ coils};

  % Coordinate system indicator (right of diagram)
  \begin{scope}[shift={(4.0,-2.0)}]
    % x-axis (left to right)
    \draw[-{Stealth[length=2mm,width=1.5mm]}] (0,0) -- (0.8,0);
    \node[right] at (0.8,0) {$x$};
    % y-axis (down to up)
    \draw[-{Stealth[length=2mm,width=1.5mm]}] (0,0) -- (0,0.8);
    \node[above] at (0,0.8) {$y$};
    % z-axis (into page)
    \fill[white] (0,0) circle (0.15);
    \node at (0,0) {$\odot$};
    \node[below left] at (0,0) {$z$};
  \end{scope}

\end{tikzpicture}
}
  \caption{Magnetic field configuration for the rubidium absorption
    cell. The large annulus represents the Helmholtz coils providing
    the primary static magnetic field. Two pairs of smaller rectangular
    coils (top/bottom and left/right) provide additional field
    control and RF magnetic fields.}
  \label{fig:helmholtz_coils}
\end{figure}

\end{document}
